\section{Sylvester's positive prime cube-sum cases}
\label{sec:sylvester}

Sylvester's positive prime cases ask whether every prime
$\ell\equiv4,7,8\pmod9$ is a sum of two rational cubes.
Dasgupta--Voight proved the $4,7$ cases when $3$ is not a cube modulo
$\ell$ \citep[Theorem~1.2.1]{DasguptaVoight2018}.
Yin's recent construction removes this restriction, with a companion
Gross--Zagier formula giving analytic rank one
\citep{Yin2026Sylvester47,Yin2026SylvesterGZ}.
Burungale--Tian prove the remaining $8$ case and state the resulting
full prime cube-sum theorem
\citep[Theorems~1.1 and~1.2]{BurungaleTian2026Sylvester}.
The ramified-prime CM converse and its Sylvester consequence are also
stated in the earlier preprint of \citet[Theorems~1.1 and~1.3]{Kriz2022}.

The proof below instead deduces these cases uniformly from
Theorem~\ref{thm:main} at $p=3$. This is a prime of additive reduction
for the cube-sum curves. The role of a rank-one converse at this prime
is also explained by
\citet[Section~1.2.2]{BurungaleTian2026Sylvester}.
Our input is Satg\'e's classical paired isogeny descent: retaining
its Tate--Shafarevich terms bounds the full Selmer corank, rather than
only the Mordell--Weil rank.

\begin{corollary}[Sylvester's positive prime cases]
\label{cor:sylvester}
Let $\ell$ be a rational prime with $\ell\equiv4,7,8\pmod 9$, and equip
the projective cubic
\[
 E_\ell:\quad X^3+Y^3=\ell Z^3
\]
with origin $(1:-1:0)$. Then
\[
 \ord_{s=1}L(E_\ell,s)=\rank_\Z E_\ell(\Q)=1,
 \qquad \#\Sha(E_\ell/\Q)<\infty.
\]
In particular, there are $x,y\in\Q$ such that $x^3+y^3=\ell$.
\end{corollary}

\begin{proof}
Use the $\Q$-isomorphic Weierstrass model
$E_\ell:y^2=x^3-27(4\ell)^2$ and its isogenous companion
$E'_\ell:y^2=x^3+(4\ell)^2$. Satg\'e's degree-three isogeny and its
dual are
\[
 \lambda:E_\ell\longrightarrow E'_\ell,\qquad
 \lambda':E'_\ell\longrightarrow E_\ell,\qquad
 \lambda'\circ\lambda=[3].
\]
These are the models and directions of Section~I, page~336, of
\citet{Satge1987CubeSums}. Let $S$ and $S'$ be the Selmer groups for
$\lambda$ and $\lambda'$, respectively, with their local Kummer
conditions. Write $\Sha(E_\ell/\Q)[\lambda]$ for the kernel of the map
on $\Sha$ induced by $\lambda$, and similarly for $\lambda'$, and put
\[
 a=\dim_{\F_3}\Sha(E_\ell/\Q)[\lambda],\qquad
 b=\dim_{\F_3}\Sha(E'_\ell/\Q)[\lambda'].
\]
These dimensions are finite: Satg\'e's descent sequences (1) and (1')
exhibit the two kernels as quotients of $S$ and $S'$.
For $D=\ell$, Proposition~1.2 of \citet{Satge1987CubeSums} gives
\[
 \dim_{\F_3}S=\dim_{\F_3}S'=1.
\]
Since $\ell>3$, his equation~(2), obtained from Lemma~1.1 and the two
descent sequences, applies and yields the exact equality
\[
 \rank_\Z E_\ell(\Q)
 =\dim_{\F_3}S+\dim_{\F_3}S'-1-a-b
 =1-a-b.
\]

The isogeny maps on Weil--Ch\^atelet groups commute with localization,
so they induce functorial maps on $\Sha$ with
$\lambda'_*\circ\lambda_*=[3]$. Thus $\lambda_*$ maps
$\Sha(E_\ell/\Q)[3]$ into $\Sha(E'_\ell/\Q)[\lambda']$, and the kernel
of this restriction is $\Sha(E_\ell/\Q)[\lambda]$. Consequently
\[
 \dim_{\F_3}\Sha(E_\ell/\Q)[3]\leq a+b.
\]
The full Kummer sequence is
\[
 0\longrightarrow E_\ell(\Q)\otimes\Q_3/\Z_3
 \longrightarrow \Sel_{3^\infty}(E_\ell/\Q)
 \longrightarrow \Sha(E_\ell/\Q)[3^\infty]
 \longrightarrow 0.
\]
The last group is a cofinitely generated $\Z_3$-module, as a quotient
of the full Selmer group. The structure theorem for its Pontryagin
dual gives
\[
 \corank_{\Z_3}\Sha(E_\ell/\Q)[3^\infty]
 \leq \dim_{\F_3}\Sha(E_\ell/\Q)[3].
\]
Combining these facts gives the required full-corank bound
\[
 s_3(E_\ell)
 =\rank_\Z E_\ell(\Q)
  +\corank_{\Z_3}\Sha(E_\ell/\Q)[3^\infty]
 \leq (1-a-b)+(a+b)=1.
\]
No finiteness of $\Sha$ has been assumed in obtaining this bound.

Now $s_3(E_\ell)\in\{0,1\}$, so Theorem~\ref{thm:main} at $p=3$
gives equality of analytic rank, Mordell--Weil rank, and this corank,
as well as finiteness of the entire group $\Sha(E_\ell/\Q)$.
For the displayed prime classes, the functional-equation sign is $-1$
by \citet[Section~1.1, equation~(1.1.3)]{DasguptaVoight2018}.
The functional equation therefore forces $L(E_\ell,1)=0$, excluding
analytic rank zero. The common rank is consequently one. This step
uses no separate parity hypothesis.

A non-torsion rational point on the projective cubic is not the origin.
The origin is its only rational point with $Z=0$, so such a point has
$Z\ne0$. Dividing the cubic equation by $Z^3$ gives the asserted
rational cube-sum representation.
\end{proof}
